% This is text associated with the open textbook "Open Electromagnetics"
% See immediately below for author, date
% <License info goes here (boilerplate, not read by project compiler)>

%ORB TITLE About This Book
%ORB AUTHOR 0        # S.W. Ellingson, Virginia Tech (ellingson@vt.edu)
%ORB DATE 20180814   # yyyymmdd
%ORB ABSTRACT Traditional preface information (i.e., what it covers, who this book is for, etc)

\label{m0146_Preface-About_This_Book} % <-- Must match name of this file and module directory name.  Don't mess with this!

{\bf Goals for this book.}
This book is intended to serve as a primary textbook for a one-semester introductory course in undergraduate engineering electromagnetics, 
including the following topics: 
electric and magnetic fields; 
electromagnetic properties of materials;
electromagnetic waves;
and devices that operate according to associated electromagnetic principles including 
resistors, 
capacitors, 
inductors, 
transformers, 
generators,
and
transmission lines.

This book employs the ``transmission lines first'' approach, in which transmission lines are introduced using a lumped-element equivalent circuit model for a differential length of transmission line, leading to one-dimensional wave equations for voltage and current.\footnote{Are you an instructor who is not a fan of the ``transmission lines first'' approach?  Then see ``{\bf What are those little numbers in square brackets?}'' later in this section.} 
This is sufficient to address transmission line concepts, including characteristic impedance, input impedance of terminated transmission lines, and impedance matching techniques.
Attention then turns to electrostatics, magnetostatics, time-varying fields, and waves, in that order.
 
{\bf What's new.}
This version of the book is the \emph{second} public release of this book.
The first release, known as ``Volume~1 (Beta),'' was released in January 2018.
Improvements from the beta version include the following:
%
\begin{itemize}
\item Correction of errors identified in the beta version errata and many minor improvements.
\item Addition of an index.
\item Accessibility features: Figures now include ``alt text'' suitable for screen reading software.
\item Addition of a separate manual of examples and solutions (see the web site).
\item Addition of source files for the book (see the web site).
%\item Figures suitable for use in presentation slides are available via the associated web site. 
\end{itemize}
 
{\bf Target audience.}
This book is intended for electrical engineering students in the third year of a bachelor of science degree program.
It is assumed that readers are familiar with the fundamentals of electric circuits and linear systems, which are normally taught in the second year of the degree program.  
It is also assumed that readers have received training in basic engineering mathematics, including 
complex numbers,
trigonometry,
vectors,
partial differential equations,
and multivariate calculus.
Review of the relevant principles is provided at various points in the text.
In a few cases (e.g., phasors, vectors) this review consists of a separate stand-alone section.

{\bf Notation, examples, and highlights.}
Section~\ref{m0005_Notation} summarizes the mathematical notation used in this book.
Examples are set apart from the main text as follows:

%%%%%%%%%%%%%%%%%%%%%%%%%%%%%%%%%%%%%%%%%%%%%%%
%ORB EXAMPLE
~\\ \begin{example} 
\begin{mdframed}[backgroundcolor=brown!20,linecolor=brown!20]\raggedright
This is an example.
%\\~\\
%Main body of example...
\end{mdframed}
% note figures will need to go between "\end{mdframed}" and "\end{example}"
\end{example}
%%%%%%%%%%%%%%%%%%%%%%%%%%%%%%%%%%%%%%%%%%%%%%%

``Highlight boxes'' are used to identify key ideas as follows:

%ORB HIGHLIGHT
\begin{mdframed}[backgroundcolor=yellow!20,nobreak=true] 
This is a key idea.
\end{mdframed}

{\bf What are those little numbers in square brackets?}
This book is a product of the 
\hyperref[m0148_Preface-About_OpenEM]{\emph{Open Electromagnetics Project}}. %({\color{red}\url{URL-TBD}}).
This project provides a large number of sections (``modules'') which are assembled (``remixed'') to create new and different versions of the book.
The text ``{\color{blue}\tiny[m0146]}'' that you see at the beginning of this section uniquely identifies the module within the larger set of modules provided by the project.
This identification is provided because different remixes of this book may exist, each consisting of a different subset and arrangement of these modules. 
%Readers can use module identifiers to find additional resources available at the \href{URL-TBD}{Open Electromagnetics website}, which is common to all remixes. 
Prospective authors can use this identification as an aid in creating their own remixes.

{\bf Why do some sections of this book seem to repeat material presented in previous sections?}
In some remixes of this book, authors might choose to eliminate or reorder modules.  
For this reason, the modules are written to ``stand alone'' as much as possible.
As a result, there may be some redundancy between sections that would not be present in a traditional (non-remixable) textbook.
While this may seem awkward to some at first, there are clear benefits: In particular, it never hurts to review relevant past material before tackling a new concept.
And, since the electronic version of this book is being offered at no cost, there is not much gained by eliminating this useful redundancy.

{\bf Why cite Wikipedia pages as additional reading?}
Many modules cite Wikipedia entries as sources of additional information.
Wikipedia represents both the best and worst that the Internet has to offer.
Most authors of traditional textbooks would agree that citing Wikipedia pages as primary sources is a bad idea, since quality is variable and
content is subject to change over time.
On the other hand, many Wikipedia pages are excellent, and serve as useful sources of relevant information that is not strictly within the scope of the curriculum.  
Furthermore, students benefit from seeing the same material presented differently, in a broader context, and with additional references cited by Wikipedia pages. 
We trust instructors and students to realize the potential pitfalls of this type of resource and to be alert for problems.

{\bf Acknowledgments.}
Here's a list of talented and helpful people who contributed to this book:\\
~\\
\emph{The staff of VT Publishing, University Libraries, Virginia Tech:}\\
Editor: Anita Walz\\
Advisors: Peter Potter, Corinne Guimont\\
Cover: Robert Browder, Anita Walz\\
~\\
\emph{Other VT contributors:}\\
Assessment: Tiffany Shoop, Anita Walz\\
Accessibility: Christa Miller\\
~\\
\emph{Virginia Tech students:}\\
Alt text writer: Stephanie Edwards\\
Figure designer: Michaela Goldammer\\
Figure designer: Kruthika Kikkeri\\
Figure designer: Youmin Qin\\
~\\
\emph{Copyediting:}\\
Melissa Ashman, Kwantlen Polytechnic University \\
~\\
\emph{External reviewers:}\\
Samir El-Ghazaly, University of Arkansas\\
Stephen Gedney, University of Colorado Denver\\
Randy Haupt, Colorado School of Mines\\
Karl Warnick, Brigham Young University\\
~\\
Also, thanks are due to the students of the Spring 2018 and Summer 2018 sections of ECE3105 at Virginia Tech who used the beta version of this book and provided useful feedback.  Thanks also to Justin Yonker, instructor of the Summer 2018 section.

